% Generated draft: requires mathematical and source review.
\documentclass[11pt]{article}
\usepackage[margin=25mm]{geometry}
\usepackage{fontspec}
\setmainfont{Latin Modern Roman}
\usepackage{amsmath,amssymb,mathtools,unicode-math}
\setmathfont{Latin Modern Math}
\usepackage{enumitem,needspace,xurl,hyperref}
\setlength{\parindent}{0pt}
\setlength{\parskip}{5pt}
\setlength{\emergencystretch}{3em}
\newcommand{\N}{\mathbb{N}}
\newcommand{\Z}{\mathbb{Z}}
\newcommand{\Q}{\mathbb{Q}}
\newcommand{\R}{\mathbb{R}}
\newcommand{\C}{\mathbb{C}}
\newcommand{\F}{\mathbb{F}}
\newcommand{\A}{\mathbb{A}}
\newcommand{\PP}{\mathbb{P}}
\newcommand{\E}{\mathbb{E}}
\newcommand{\eps}{\varepsilon}
\newcommand{\dd}{\,\mathrm{d}}
\begin{document}
\section*{UnsolvedMath ID: 30005376}
\textbf{Source problem code:} catalogue entry
\subsection{Definitions and mathematical statement}
Let \(X=\mathbb{CP}^{2}\#\mathbb{CP}^{2}\) be the oriented connected sum of two complex projective planes with their complex orientations. Let \(H_1,H_2\in H_2(X;\Z)\) be the classes of projective lines in the respective summands, so that \(H_1^2=H_2^2=1\) and \(H_1\cdot H_2=0\). For a class \(a\) admitting a connected, closed, oriented, smoothly embedded surface representative, define
\[
 g_X(a)=\min\{g(\Sigma):\Sigma\subset X\text{ is connected, closed, oriented, and smoothly embedded, with }[\Sigma]=a\}.
\]
Here genus refers to the genus of an oriented real surface.

For positive integers \(n,m\), choose smooth complex plane curves \(C_n,C_m\subset\mathbb{CP}^{2}\) of degrees \(n,m\). These curves have complex dimension one and real dimension two; this is the meaning of the source's term ``complex surfaces'' in this passage. Form their oriented connected sum \(\Sigma_{n,m}=C_n\# C_m\subset X\), using the two respective ambient summands. Then
\[
 [\Sigma_{n,m}]=nH_1+mH_2,\qquad
 g(\Sigma_{n,m})=\frac{(n-1)(n-2)}{2}+\frac{(m-1)(m-2)}{2}.
\]
\textbf{General question (Question 9 in the supplied report passage).} Is it true that, for every pair \((n,m)\in\Z_{>0}^{2}\),
\[
 g_X(nH_1+mH_2)=\frac{(n-1)(n-2)}{2}+\frac{(m-1)(m-2)}{2}?
\]
Equivalently, does every connected, closed, oriented, smoothly embedded surface \(\Sigma\subset X\) representing \(nH_1+mH_2\) satisfy
\[
 g(\Sigma)\geq\frac{(n-1)(n-2)}{2}+\frac{(m-1)(m-2)}{2}?
\]

\textbf{Specific question retained from the source (Question 8).} For every integer \(k>1\), is
\[
 g_X(2kH_1+6kH_2)
 =\frac{(2k-1)(2k-2)}{2}+\frac{(6k-1)(6k-2)}{2}
 =20k^2-12k+2?
\]
The variable \(k\) here denotes the variable called \(n\) in Question 8 of the source.

\textbf{Source qualifications.} The supplied dated curation labels the record ``partially solved'' and attributes to Ait Nouh infinite counterexample families disproving Lawson's general minimum-genus formula. It also states that the degree-\(2k\), degree-\(6k\) subfamily and related positive cases are not completely resolved. Thus the general question above is retained as the source's mathematical target, without asserting that its universal affirmative formulation remains open. These literature claims have not been independently verified for this draft.

The report passage states that Question 9 has an affirmative answer for some specific small positive degrees, without identifying those pairs here. It also states that, for every integer \(d>2\), a degree-\(d\) complex plane curve included in one summand fails to minimize genus in the class \(dH_1=(d,0)\); this concerns a zero second coordinate, outside the positive-degree domain of Question 9. The passage further states that its Theorem 5 answers its Question 7 affirmatively for parameter \(n=1\), and that an affirmative answer to Question 8 would imply a negative answer to Question 7. Neither Question 7 nor Theorem 5 is supplied, so these contextual assertions cannot be made into self-contained mathematical statements here. The report presents Question 9 as motivated by the Thom Conjecture, citing its reference [9], and refers to work of Nouh in its reference [11]; the complete entries for those numbered references are absent from the supplied record.
\subsection{Short English statement}
The question asks whether joining smooth complex plane curves of positive degrees \(n\) and \(m\) gives the smallest possible genus among connected, closed, oriented smooth surfaces representing \(nH_1+mH_2\) in \(\mathbb{CP}^{2}\#\mathbb{CP}^{2}\). The source also singles out degrees \(2k\) and \(6k\), with \(k>1\). Supplied dated curation reports counterexamples to the general formula and incomplete resolution of the narrower family; this draft does not independently verify those assessments.
\subsection{Sources}
\begin{itemize}
\item Local UnsolvedMath registry, permanent ID \(30005376\), title ``Minimum Genus Connected Sums of Complex Plane Curves,'' exact source problem code \texttt{catalogue entry}, in the Oberwolfach Reports Open Problems collection. No individual proposer is supplied; the proposed year is 2023.
\item Supplied report citation: \emph{Low-dimensional Topology} (2023), Oberwolfach Reports, pp.\ 151--213, DOI 10.4171/owr/2023/3. \url{https://doi.org/10.4171/owr/2023/3}. The supplied excerpt contains Questions 8 and 9 and contextual references to Question 7, Theorem 5, and references [9] and [11]. Oberwolfach Reports are described in the record as documenting mathematical workshops and their open-problem discussions.
\item Supplied literature reference, labeled ``Ait Nouh, minimal genus in CP2 connected sum CP2'': \url{https://doi.org/10.2140/agt.2014.14.671}. This is the supplied attribution for the curation's counterexample claim; a complete bibliographic entry is not supplied.
\item Supplied literature review, checked \texttt{2026-08-22}, lists both \url{https://doi.org/10.2140/agt.2014.14.671} and \url{https://doi.org/10.4171/owr/2023/3}. Its ``partially solved'' label and assessment concerning Lawson's general formula and the narrower subfamily are supplied dated curation, not independently verified literature findings. The background explicitly requests independent verification of time-sensitive status and cited claims.
\item The registry labels the statement \texttt{corrected\_verified} and records that the official report was checked and the notation repaired for self-contained publication. This is supplied source-verification metadata, not a claim that this draft's preparer visited or verified either external source.
\end{itemize}
\end{document}
